\documentclass{article}



% Language setting

% Replace `english' with e.g. `spanish' to change the document language

\usepackage[english]{babel}



% Set page size and margins

% Replace `letterpaper' with `a4paper' for UK/EU standard size

\usepackage[letterpaper,top=2cm,bottom=2cm,left=3cm,right=3cm,marginparwidth=1.75cm]{geometry}



% Useful packages

\usepackage{amsmath}

\usepackage{amssymb}

\usepackage{graphicx}

\usepackage[colorlinks=true, allcolors=black]{hyperref}

\usepackage{titlesec}
\usepackage{xparse}



\NewDocumentCommand{\qcq}{m o o m}{%
	\ensuremath{\forall #1 \IfValueT{#2}{, #2} \IfValueT{#3}{, #3}  \in  \mathbb{#4}}%
}
\NewDocumentCommand{\ext}{m o o m}{%
	\ensuremath{\exists #1 \IfValueT{#2}{, #2} \IfValueT{#3}{, #3}  \in  \mathbb{#4}}%
}



\title{\textbf{Exercices sur les fonctions logarithmiques et exponentiels}}

\author{\textbf{Yassin Akkab}}



\newcommand{\R}{\mathbb{R}}
\newcommand{\N}{\mathbb{N}}
\newcommand{\Z}{\mathbb{Z}}
\newcommand{\C}{\mathbb{C}}
\newcommand{\Q}{\mathbb{Q}}
\newcommand{\D}{\mathbb{D}}





\graphicspath{ {./Screenshots/} }




\begin{document}
	
	\maketitle
	
	\subsection*{Exercice 1 : Étude de la fonction exponentielle}
	Étudiez la continuité et la dérivabilité de la fonction \( f(x) = e^x \) sur \( \mathbb{R} \).
	\begin{itemize}
	
		\item Calculez la limite de \( e^x \) lorsque \( x \to +\infty \) et \( x \to -\infty \).
			\item Déterminez la dérivée de \( f(x) \).
	\end{itemize}
	
	\subsection*{Exercice 2 : Étude de la fonction logarithme népérien}
	Étudiez la continuité et la dérivabilité de la fonction \( g(x) = \ln(x) \) sur \( ]0, +\infty[ \).
	\begin{itemize}
	
		\item Calculez la limite de \( \ln(x) \) lorsque \( x \to 0^+ \) et \( x \to +\infty \).
			\item Déterminez la dérivée de \( g(x) \).
	\end{itemize}
	
	\subsection*{Exercice 3 : Limites et dérivées combinées}
	Calculez les limites suivantes :
	\[
	\lim_{x \to 0^+} \frac{\ln(x)}{x}, \quad \lim_{x \to 0^+} x \ln(x), \quad \lim_{x \to 0^+} \frac{1 - e^x}{x}
	\]
	

	
	\subsection*{Exercice 4 : Étude de la fonction \( f(x) = e^{-x} \)}
	Étudiez la continuité et la dérivabilité de \( f(x) = e^{-x} \) sur \( \mathbb{R} \).
	\begin{itemize}
		\item Déterminez la dérivée de \( f(x) \).
		\item Calculez les limites de \( e^{-x} \) lorsque \( x \to +\infty \) et \( x \to -\infty \).
	\end{itemize}
	
	\subsection*{Exercice 5 : Dérivée de la fonction composée}
		Etudiez les fonctions suivantes : 
	\[
	f(x) = e^{x^2}, \quad g(x) = \ln(x^2 + 1)
	\]

	
	
	
	
	
	
	
	
	
	
	
	
\end{document}